\section{Загальний опис системи}
\label{sec:zahalnyi-opys}

\subsection{Загальна характеристика та перспектива продукту}

Система є самостійним веб-застосунком без інтеграцій із зовнішніми сервісами в межах MVP. Кожен зареєстрований користувач має власний ізольований простір даних~--- каталог, добірки та статистику. Система не є частиною більшої платформи і не надає публічного API у версії~v1. \sked{K38,K41,K42}

\subsection{Склад функціональних блоків}

Система реалізує такі функціональні блоки:
\begin{enumerate}
  \item Автентифікація та управління акаунтом.
  \item Особистий каталог медіа-контенту.
  \item Управління елементами контенту (CRUD).
  \item Теги та фільтрування за тегами.
  \item Sessions~--- фіксація повторних переглядів або прочитань.
  \item Добірки~--- власні іменовані групи елементів.
  \item Автозаповнення дати завершення.
  \item Особиста статистика.
\end{enumerate}
\sked{K10,K13,K16,K18,K21,K27,K28}

\subsection{Ролі користувачів}

У версії~v1 існує єдина роль. Адміністративна роль відсутня.

\begin{table}[h]
\caption{Ролі користувачів системи}
\label{tab:roli}
\begin{tabular}{L{3cm} L{10cm}}
\toprule
\textbf{Роль} & \textbf{Опис} \\
\midrule
Користувач (User) & Зареєстрований власник акаунту. Має доступ виключно до власного каталогу, добірок і статистики. Може виконувати всі CRUD-операції над власними даними. \\
\bottomrule
\end{tabular}
\end{table}
\sked{K44}

\subsection{Операційне середовище}

Система реалізується як веб-застосунок. Нативний мобільний застосунок у версії~v1 не передбачений. Система повинна коректно функціонувати в останніх двох версіях браузерів Chrome, Firefox, Edge та Safari, включно з мобільним Safari. Мова реалізації бекенду не регламентується цим документом і визначається на етапі проектування архітектури. \sked{K38,K39,K45}

\subsection{Обмеження проектування та реалізації}

\begin{itemize}
  \item Реалізація виключно у вигляді веб-застосунку.
  \item У версії~v1 підтримується лише один мовний інтерфейс.
  \item Система не інтегрується із зовнішніми базами медіа-контенту (TMDB, Google Books тощо).
  \item Кожен користувач має доступ виключно до власних даних; крос-користувацький доступ не передбачений.
\end{itemize}
\sked{K38,K40,K41,K42}

\subsection{Межі MVP~--- що не входить до версії~v1}

Наступні функції є технічно можливими, але свідомо виключені з MVP для стримання обсягу першої версії. Детальний перелік наведено у Додатку~А. \sked{K43,K50}
